\documentclass[11pt,a4paper]{article}
\usepackage[utf8]{inputenc}
\usepackage[T1]{fontenc}
\usepackage[english]{babel}
\usepackage{lmodern,microtype,amsmath,amssymb,booktabs,geometry,xcolor,fancyhdr}
\usepackage[hidelinks]{hyperref}
\geometry{margin=2.3cm,headheight=14pt}
\setlength{\parindent}{0pt}
\setlength{\parskip}{0.4em}
\setlength{\emergencystretch}{3em}
\definecolor{gcblue}{RGB}{34,73,110}
\pagestyle{fancy}\fancyhf{}
\fancyhead[L]{\small GC2D: shared extended-space implementation}
\fancyhead[R]{\small Numerical architecture}\fancyfoot[C]{\thepage}
\newcommand{\R}{\mathbb{R}}
\begin{document}
\begin{center}
{\LARGE\bfseries\color{gcblue}ABBA and BM4 in Extended Space}\par
\vspace{0.35em}
{\large Common composition, spatial projection and passive energy}
\end{center}

\section{Scope and public methods}
The directory \path{src/methods/extended/} contains the common
implementation of the existing ABBA and BM4 methods. It introduces no public
method named Extended. Public class names and exports from \texttt{simulation} are retained. ABBA6
now uses the same single outer projection as ABBA4. Its observed substeps
become unprojected pairs and its work arrays have one column per step.
Internal imports use the defining modules. Old ABBA and BM4 import adapters
and historical \texttt{simulation.methods} routes have been removed.

For $N$ independent planar particles, the physical state $z\in\R^{2N}$ is
packed in component-major order. The splitting acts on two copies
$Y=(u,v)\in\R^{4N}$. Accepted copies coincide. Optional energy tracking adds
$N$ time entries and $N$ physical conjugate momenta $\kappa$, giving a
$6N$-component accepted internal workspace. These auxiliary variables never
enter a spatial projection equation. Physical output remains in $\R^{2N}$.

\section{Direct and adjoint compositions}
Let $\chi_s$ denote a first-order map on the extended spatial workspace, with
adjoint $\chi_s^*=\chi_{-s}^{-1}$, including the appropriate time translation
for nonautonomous dynamics. A recipe with $2r$ stages is
\[
 \Psi_h=\chi_{b_{2r}h}\circ\chi^*_{b_{2r-1}h}\circ\cdots
          \circ\chi_{b_2h}\circ\chi^*_{b_1h}.
\]
In execution order the stages alternate adjoint, direct. The weights obey
$b_j=b_{2r+1-j}$ and $\sum_j b_j=1$. Reflection exchanges each map with its
adjoint. These properties are checked by the recipe constructor; additional
order conditions remain necessary for higher-order accuracy.

With stage clock $\tau$ and signed duration $s=b_jh$, the adjoint evaluates at
$\tau$, the direct map at $\tau+s$, and the clock subsequently becomes
$\tau+s$. Negative coefficients reverse the clock as well as the subflow.
For uncoupled GC maps, the pair with weights $(1/2,1/2)$ reproduces the
endpoint-time A--B--B--A construction. Each direct or adjoint map contains two
spatial shears. BM4 can additionally apply the existing exact harmonic
coupling of the two spatial copies, in reversed operation order for the
adjoint.

ABBA4 uses three unprojected ABBA pairs scaled by
\[
 \gamma=\frac{1}{2-2^{1/3}},\qquad
 \delta=-\frac{2^{1/3}}{2-2^{1/3}},\qquad
 (\gamma,\delta,\gamma).
\]
Its six direct/adjoint weights are
$(\gamma/2,\gamma/2,\delta/2,\delta/2,\gamma/2,\gamma/2)$.
BM4 retains its twelve mirrored weights and its own fourth-order recipe.
ABBA6 uses seven \emph{unprojected} ABBA2 pairs with the Yoshida weights.
Each weight is repeated twice after division by two, producing fourteen
direct/adjoint stages. One symmetric projection encloses the complete
composition. This deliberately replaces the former seven-projection map.

\clearpage
\section{One common spatial projection equation}
Write $m=2N$, $Ez=(z,z)$, $G=[I,-I]$ and $N_0=G^{\mathsf T}$. The reduced
Hairer projection finds $\mu\in\R^m$ such that
\[
 r(\mu)=G\Psi_h(Ez+N_0\mu)+2\mu=0.
\]
After convergence the two corrected outputs are averaged. With
$P=\tfrac12[I,I]$, the accepted state is
\[
 z_+=P\bigl(\Psi_h(Ez+N_0\mu)+N_0\mu\bigr).
\]
The same equation applies to ABBA2, the complete ABBA4 and ABBA6
compositions, and BM4; only $\Psi_h$ changes. All have one nonlinear solve
per complete step.

The simultaneous ABBA formulation solves for $(u_+,v_+,\mu)\in\R^{3m}$:
\[
 \begin{pmatrix}
 u_+-\mu-[\Psi_h(Ez+N_0\mu)]_u\\
 v_++\mu-[\Psi_h(Ez+N_0\mu)]_v\\
 u_+-v_+
 \end{pmatrix}=0.
\]
Arithmetic midpoint variants instead start from $Ez$, apply the complete
recipe, and average once. They do not solve a nonlinear equation; arithmetic
projection does not imply exact symplecticity or reversibility.

\section{Nonlinear solution and exact tangents}
All spatial roots use the unchanged stopping threshold
\[
 \mathrm{atol}+\mathrm{rtol}\max(1,\|z\|_\infty).
\]
Newton and good Broyden share the convergence machinery in
\path{src/methods/_nonlinear.py}. Broyden's initial reduced
Jacobian is $4I$, the exact zero-step value of $Dr$. The simultaneous
zero-step block matrix is
\[
 \begin{pmatrix} I_{2m}&-2N_0\\G&0\end{pmatrix}.
\]
Broyden subsequently uses residual secants and evaluates no analytic
Jacobian.

If $M=D\Psi_h$, Newton's reduced Jacobian is $GMN_0+2I$. The common
composition retains the input to each spatial shear. Exact stage tangents
are evaluated at those inputs and multiplied in execution order,
$M=M_{2r}\cdots M_1$, without replaying vector-field stages.
Independence of the particles gives $4\times4$ blocks for $M$, $2\times2$
reduced Newton blocks, and $6\times6$ simultaneous blocks. BM4 also retains
its optional dense centered-difference differentiation of the complete map.

\clearpage
\section{Passive energy from accepted stages}
The dynamics supply $H(t,z)$ and $-\partial_tH(t,z)$. For the accepted spatial
trace, let $(t_j,s_j,z_j)$ denote the evaluation time, signed duration and
physical source state of each individual shear. The shared quadrature is
\[
 \Delta\kappa=\frac12\sum_j s_j\,[-\partial_tH(t_j,z_j)].
\]
There are four shear points for ABBA2, twelve for ABBA4, 28 for ABBA6,
and 24 for BM4. The factor $1/2$ converts summed doubled momentum to the
physical variable $\kappa$. Spatial harmonic coupling does not change this
passive accumulator. Energy derivatives are evaluated only after the spatial
root converges, from its retained trace; trial roots do not evaluate energy.
The accepted diagnostic is
\[
 H(t_n,z_n)+\kappa_n-H(t_0,z_0),\qquad \kappa_0=0.
\]
This is an energy-balance diagnostic, not a claim that the time-dependent
physical Hamiltonian is conserved. Tracking does not modify spatial states,
nonlinear unknown dimensions or stopping scales.

\section{Execution and organization}
Initialization binds the method's recipe and projection options once. The
ordinary implicit execution is
\[
 \texttt{advance}\;\longrightarrow\;\texttt{solve\_projection}
 \;\longrightarrow\;\texttt{compose}
 \;\longrightarrow\;\texttt{direct\_map / adjoint\_map}.
\]
ABBA4 and ABBA6 share this execution exactly; only their recipes differ.
Public configurations live in \texttt{abba.py} and \texttt{bm4.py}; common
numerical operations live in \texttt{core/}, and optional event conversion
lives in \texttt{observations.py}. The general
integration coordinator remains responsible for main steps, saved times,
shadow output samples, collection and observer dispatch. Observers receive
independent snapshots in the existing public event types; BM4 events reuse
accepted traces without repeating the spatial composition.

Current diagrams and module responsibilities are documented in
\path{docs/models/extended/simulation/extended-simulation-architecture.md}.
Tests verify public identity, coefficient and projection controls, signed
nonautonomous maps, exact tangents and passive tracking. ABBA6 tests also
verify sixth-order convergence on a nonlinear nonautonomous field with both
root formulations and solvers. Historical ABBA6 results with seven separate
projections represent a different map and must be recomputed for comparisons.
\end{document}
